\documentclass[11pt,a4paper]{article}

% Report template for the truncated SHA-2 assignment.
% Replace each gray italic prompt before you submit.
% Build the table of contents with two passes, or with latexmk:
%   latexmk -pdf report.tex

\usepackage[T1]{fontenc}
\usepackage[utf8]{inputenc}
\usepackage{lmodern}
\usepackage[margin=2.5cm]{geometry}
\usepackage{microtype}
\usepackage{setspace}
\usepackage{amsmath,amssymb}
\usepackage{booktabs}
\usepackage{array}
\usepackage{enumitem}
\usepackage{xcolor}
\usepackage[hidelinks]{hyperref}

\onehalfspacing
\setlist[itemize]{topsep=0.3em,itemsep=0.15em}
\setlist[enumerate]{topsep=0.3em,itemsep=0.2em}

% Edit these four lines.
\newcommand{\studentname}{Your Name}
\newcommand{\studentid}{Student ID}
\newcommand{\coursename}{Data Security}
\newcommand{\reporttitle}{How Secure Is a Truncated SHA-2 Digest?}

\newcommand{\Hn}{H_{n}}
\newcommand{\prefix}{\operatorname{prefix}}
\newcommand{\prompt}[1]{%
  \par\smallskip
  {\color{black!45}\itshape #1\par}%
  \smallskip
}

\hypersetup{
  pdftitle={\reporttitle},
  pdfauthor={\studentname}
}

\begin{document}
% Alphabetic anchors keep the title page from sharing page.1 with the body.
\pagenumbering{alph}

\begin{titlepage}
  \centering
  \thispagestyle{empty}
  \vspace*{1.6cm}
  {\LARGE\bfseries \reporttitle\par}
  \vspace{0.8cm}
  {\Large Course report\par}
  \vspace{1.6cm}
  {\large Truncated SHA-256, birthday collisions,\\
   and a signature that checks only a prefix\par}
  \vfill
  \begin{tabular}{@{}>{\bfseries}ll@{}}
    Student  & \studentname \\
    ID       & \studentid \\
    Course   & \coursename \\
    Date     & \today \\
  \end{tabular}
  \vfill
  {\small Problem statement: \texttt{assignment/problem/problem.tex}\par}
  \vspace*{1cm}
\end{titlepage}

\pagenumbering{roman}
\tableofcontents
\clearpage
\pagenumbering{arabic}

\section{Introduction}
\prompt{%
  State the question in your own words.
  A signature that verifies \(\Hn(x) = \prefix_{n}(H(x))\) accepts every
  message with that same \(n\)-bit prefix.
  Say which \(n\) you studied and what you concluded about the signature.%
}

\section{Background}

\subsection{SHA-2}
\prompt{%
  Identify the member of the family you hashed (SHA-256 unless you changed
  it), its output length, and the properties the full function is designed
  to provide: collision resistance and preimage resistance.
  Name where a digest of this kind is used (signatures, TLS, blockchains)
  only as far as it matters for the attack.%
}

\subsection{Truncation and significant bits}
\prompt{%
  Define \(\Hn(x) = \prefix_{n}(H(x))\) and state that the prefix is the
  most significant \(n\) bits, the leading bits when the digest is written
  in hex.
  A short example, even one byte, is enough to fix MSB and LSB.%
}

\subsection{Birthday bound}
\prompt{%
  Give the collision condition \(x \neq x'\) and \(\Hn(x) = \Hn(x')\),
  and the generic cost \(O(2^{n/2})\).
  Record the lecture's reading of that cost: truncation leaves at most
  \(n/2\) bits of collision security, short prefixes such as \(n = 32\)
  and \(n = 64\) are practical, and \(n \ge 128\) is the example of a
  length that is still large.%
}

\subsection{The signature on a prefix}
\prompt{%
  Say what the verifier computes and what it compares.
  Then say what follows when two messages share a prefix:
  a signature on one verifies on the other.%
}

\section{Alice's investigation}
\prompt{%
  Summarize the poisoned-message setting you are replaying.
  Two documents that a reader accepts as different messages share one
  digest, so the signature transfers.
  Cite Daum and Lucks, and point at the published MD5 pair only as the
  historical case.
  Keep this section to the analogy.
  The measurements belong in Section~\ref{sec:results}.%
}

\section{Collision search}
\label{sec:search}
\prompt{%
  Open with the search problem: sample inputs until two of them share an
  \(n\)-bit prefix.
  Then one subsection per approach you actually ran.
  Delete the subsections you did not use, and add one for each further
  idea (distinguished points, a different cycle finder, a GPU search,
  and so on).%
}

\subsection{Hash table}
\prompt{%
  Describe the key you store, the value you store beside it, and the test
  that declares a collision.
  Give the memory per entry and the \(n\) at which the table stopped
  fitting.
  This is the memory-heavy end of the tradeoff.%
}

\subsection{Floyd's cycle-finding algorithm}
\prompt{%
  Describe the function you iterate, the two pointers, and the extra work
  you do after a meeting point to recover a colliding pair.
  State the memory, which should stay essentially constant, and the number
  of hashes relative to \(2^{n/2}\).%
}

\subsection{Time--memory tradeoff}
\prompt{%
  Put the two previous methods on one scale.
  If you used distinguished points, chains, or another intermediate
  scheme, specify the parameter that moves memory in one direction and
  hashing work in the other, and report the settings you tried.%
}

\subsection{Parallel search}
\prompt{%
  Describe how the search was split (processes, threads, or machines),
  how a worker reported a hit, and how the speed compared with one worker.
  Note any duplicated work.%
}

\section{Experiments}
\label{sec:results}

\subsection{Setup}
\prompt{%
  Machine, core count, RAM, language, and hash implementation.
  State how you timed a run and how you measured memory.
  Fix the input format (counter, random block, chosen prefix plus a
  free block) so a reader can reproduce a pair.%
}

\subsection{Results}
\prompt{%
  Fill Table~\ref{tab:results}.
  One row per \((n, \text{method})\) you finished.
  The largest \(n\) is the bottom line of the assignment; start at 32
  and keep every successful length, including the ones that were easy.%
}

\begin{table}[ht]
  \centering
  \caption{Collision searches.
           Replace the dashes.
           ``Hashes'' is the number of SHA-256 evaluations.}
  \label{tab:results}
  \begin{tabular}{@{}clrrll@{}}
    \toprule
    \(n\) & Method & Hashes & Time & Memory & Prefix \\
    \midrule
    32 & hash table & --- & --- & --- & \texttt{--------} \\
    32 & Floyd      & --- & --- & --- & \texttt{--------} \\
    64 &            & --- & --- & --- & \texttt{----------------} \\
    \midrule
    \multicolumn{6}{@{}l}{%
      Largest \(n\) reached: \rule{2.2cm}{0.4pt}\quad
      method: \rule{3.2cm}{0.4pt}} \\
    \bottomrule
  \end{tabular}
\end{table}

\subsection{What limited the next prefix length}
\prompt{%
  For the smallest \(n\) you did not finish, say whether you ran out of
  memory, time, or both, and which method came closest.
  Compare that wall with \(2^{n/2}\).%
}

\section{Meaningful messages}
\prompt{%
  Leave this section until the lecture that builds a collision into a
  document a person will actually read.
  Then describe the two messages, where the free bytes sit, and why a
  reader sees two different texts while \(\Hn\) stays the same.
  Attach the files or quote only enough of them to identify the pair.
  If that lecture has not happened yet, write one sentence saying so
  and point ahead to the conclusion you can already draw from
  Section~\ref{sec:results}.%
}

\section{Security conclusion}
\prompt{%
  Answer the mission directly.
  For each \(n\) you broke, the signature on \(\Hn\) does not separate
  the two inputs you found.
  For the lengths you did not break, give the birthday cost and say
  whether that cost is still a practical attack or only a theoretical one.
  Use the lecture's markers: \(n = 32\) and \(n = 64\) are the short
  prefixes, and \(n \ge 128\) is the example of a prefix that remains
  large.
  Close with why the birthday bound, rather than a dedicated SHA-256
  break, is the attack that decides this question.%
}

\section{Conclusion}
\prompt{%
  Three or four sentences.
  Largest \(n\), the method that reached it, the memory--time choice
  that mattered, and the consequence for the signature.
  State the meaningful-message attack as done, or as waiting on the
  later lecture.%
}

\appendix
\section{Colliding inputs}
\label{app:pairs}
\prompt{%
  For every row of Table~\ref{tab:results}, list the two inputs and the
  shared prefix.
  Prefer a path to a file in the repository over a paste of long inputs.%
}

\begin{verbatim}
n = 32
method:
x  =
x' =
Hn =
\end{verbatim}

\addcontentsline{toc}{section}{References}
\begin{thebibliography}{9}

\bibitem{daum-lucks}
Magnus Daum and Stefan Lucks.
\newblock Hash Collisions (The Poisoned Message Attack).
\newblock \url{https://www.sos.cs.ru.nl/applications/courses/security2015/md5collisions/index.html}.

\bibitem{fips180-4}
National Institute of Standards and Technology.
\newblock \emph{Secure Hash Standard (SHS)}, FIPS PUB 180-4, 2015.

% Add a bibitem for each paper or tool you actually used, for example
% the hash library and any write-up of Floyd's algorithm you followed.

\end{thebibliography}

\end{document}
